\section{Arithmetic, walks, and entropy}\label{sec:preliminaries}

We record the conventions needed to sample a fixed walk and to compare
the resulting residue distributions. All logarithms are natural.
The notation \(A\ll_D B\) means \(A\le c_D B\) with a constant depending
only on \(D\); \(A\asymp_D B\) means both inequalities.
Constants without a subscript are absolute. Asymptotic statements about
a collection of walks or starting times will always be uniform over
that collection.

\subsection{Split primes}

The ring \(\Z[i]\) is Euclidean and hence has unique factorization.
For each rational prime \(p\equiv1\pmod4\), its two conjugate Gaussian
prime factors are nonassociate, have norm \(p\), and have residue fields
of size \(p\); see~\cite[Corollary 7.9 and Theorems 7.14, 9.7, 9.9]{Conrad}.
We refer to the choice of one of these two factors as a \emph{sign}.
Signs of distinct rational primes will be independent and fair.

We use the following consequences repeatedly. If \(\pi\mid v\) and
\(\overline\pi\mid v\), then \(p\mid v\), meaning that \(p\) divides both
integer coordinates of \(v\). If \(a\in\Z\), then
\(\pi\mid a\) if and only if \(p\mid a\).
Every nonzero multiple of \(\pi\) has length at least \(\sqrt p\).
Finally, multiplication by a Gaussian integer \(\alpha\) acts on
\(\Z^2\) by an integer matrix of determinant \(N(\alpha)\). Thus if
\(\alpha\mid v,w\), then
\[
\det(v,w)\in N(\alpha)\Z.
\]

For \(T>1\), let
\[
\mathcal B_T=\{p\text{ rational prime}:T\le p\le2T,\ p\equiv1\pmod4\},
\qquad k_T=|\mathcal B_T|,
\qquad L_T^*=\frac1{k_T}\sum_{p\in\mathcal B_T}\log p.
\]
All uses have \(T\) sufficiently large that \(k_T>0\).
For one fixed batch we abbreviate \(k_T,L_T^*\) by \(k,L_*\).
The prime number theorem in arithmetic progressions for the fixed
modulus \(4\) gives~\cite[Equations (1.1)--(1.2), p.~66]{Selberg}
\begin{equation}\label{eq:prime-count}
k_T=(1/2+o(1))\frac{T}{\log T},
\qquad \log T\le L_T^*\le\log(2T).
\end{equation}
Indeed Selberg's weighted asymptotic
\(\sum_{p\le x,\ p\equiv1(4)}\log p\sim x/2\), subtracted at \(T\)
and \(2T\), gives a sum over \(T<p\le2T\). Adding the possible
endpoint term at \(p=T\) changes it by at most \(\log T\). The weighted
sum over \(\mathcal B_T\) therefore equals \((1/2+o(1))T\) and lies
between \(k_T\log T\) and \(k_T\log(2T)\), whose ratio tends to \(1\).
Only fixed-modulus, fixed-ratio intervals are involved. In particular,
all constants implicit in \eqref{eq:prime-count} hold for every
sufficiently large \(T\).

\subsection{Entropy conventions}

For a finite variable \(X\) with law \(p\),
\(H(X)=-\sum_xp(x)\log p(x)\), with \(0\log0=0\).
We use conditional entropy and conditional mutual information in
their usual finite-alphabet forms:
\[
I(X;Y\mid Z)=H(X\mid Z)-H(X\mid Y,Z).
\]
Standard identities and inequalities may be found in
Cover and Thomas~\cite[Chapter 2]{CoverThomas}.
When residue maps are selected randomly, we first fix those maps
and compute entropy in the walk-sampling experiment; only afterward
do we average over the maps.
An expectation outside an entropy in this usage is an average of
entropies, not the entropy of a mixture that includes the map choices.

For laws \(p,q\) on a finite set, put
\(\TV(p,q)=\frac12\sum_x|p(x)-q(x)|\).
Relative entropy is
\(D_{\mathrm{KL}}(p\Vert q)=\sum_xp(x)\log(p(x)/q(x))\),
with the usual value \(+\infty\) if \(q\) vanishes where \(p>0\).
In particular, on an alphabet of size \(a\),
\[
D_{\mathrm{KL}}(p\Vert\mathrm{Unif})=\log a-H(p),
\qquad D_{\mathrm{KL}}(p\Vert q)\ge2\TV(p,q)^2.
\]
The second inequality is Pinsker's inequality in natural logarithms
\cite[Lemma 11.6.1]{CoverThomas}.

For completeness, let $P,Q$ be finite probability laws and put
\(B=\{x:P(x)>Q(x)\}\), \(u=P(B)\), and \(v=Q(B)\).
The log-sum inequality, from convexity of \(t\log t\), gives
\[
 D_{\mathrm{KL}}(P\Vert Q)
 \ge D_{\mathrm{KL}}(\operatorname{Ber}(u)\Vert
                          \operatorname{Ber}(v)).
\]
Here $\operatorname{Ber}(u)$ is the Bernoulli law with success probability $u$.
As a function of \(u\), the expression on the right has value and
first derivative zero at \(u=v\), and second derivative
\(1/[u(1-u)]\ge4\). It is therefore at least
\(2(u-v)^2=2\TV(P,Q)^2\), with boundary cases by continuity.

We shall also use \(h(u)=-u\log u-(1-u)\log(1-u)\) for binary entropy,
with its continuous values at \(0\) and \(1\).

\subsection{Sampling a deterministic walk}

Fix \(D\ge1\) and an infinite deterministic sequence
\((z_t)_{t\ge0}\) of distinct points of \(\Z[i]\), with
\(|z_{t+1}-z_t|\le D\).
A random time is a finitely supported nonnegative integer-valued
variable. All later sampling procedures produce such times.
The walk itself is never assigned a probability law.

Recall the increment bound
\[
K_D=|\{u\in\Z[i]:|u|\le D\}|.
\]
This includes the zero vector; although a self-avoiding walk never
uses it as an increment, this harmless upper bound simplifies notation.

\begin{lemma}[Displacement entropy]\label{lem:displacement}
If \(T_0,T_1\) are jointly distributed finite random times and
\(|T_1-T_0|\le n\) almost surely, then
\[
H(z_{T_1}-z_{T_0})\le2\log(n+1)+O_D(1).
\]
The same bound holds conditionally on any event or variable under
which the time-difference bound still holds.
If \(F\) is any finite vector of residue maps, then
\[
H(F(z_{T_1}))\ge H(F(z_{T_0}))-H(z_{T_1}-z_{T_0}).
\]
\end{lemma}
\begin{proof}
The displacement lies in a disk of radius \(Dn\), which contains at
most \(C_D(n+1)^2\) lattice points. Entropy is at most the logarithm
of support size. For the last assertion, \(F(z_{T_0})\) is determined
by \(F(z_{T_1})\) and the displacement. The chain rule and
subadditivity give the claimed inequality.
\end{proof}

\subsection{Continuity and short lists}

\begin{lemma}[Continuity]\label{lem:continuity}
If two laws of \((X,Y)\), on the same finite product alphabet, have
total variation at most \(\epsilon\le1/2\), and the product alphabet
has at most \(a\) elements, then
\[
|H(X\mid Y)-H(\widetilde X\mid\widetilde Y)|
\le 2h(\epsilon)+2\epsilon\log a.
\]
\end{lemma}
\begin{proof}
A coupling with mismatch probability equal to total variation yields
\(|H(U)-H(\widetilde U)|\le h(\epsilon)+\epsilon\log a\)
for variables on an alphabet of size at most \(a\).
Apply this once to the joint variables and once to their
\(Y\)-marginals, and subtract.
\end{proof}

The next lemma adapts the error-indicator proof of Fano's inequality
to lists; see~\cite[Theorem 2.10.1]{CoverThomas}.

\begin{lemma}[A short-list entropy bound]\label{lem:list}
Let \(X\) have at most \(a\) values, let \(Y\) be finite data, and let
\(\mathcal L(Y)\) be a list of possible values of \(X\).
If
\[
\Pr\{X\in\mathcal L(Y),\ |\mathcal L(Y)|\le b\}=\rho,
\qquad 1\le b\le a,
\]
then
\[
\log a-H(X\mid Y)\ge \rho\log(a/b)-h(\rho).
\]
\end{lemma}
\begin{proof}
Let $B$ indicate the event in the statement. It is determined by $(X,Y)$,
so the chain rule gives
\[
H(X\mid Y)=H(B\mid Y)+H(X\mid B,Y)
\le h(\rho)+\rho\log b+(1-\rho)\log a.
\]
On $B=1$ the conditional support lies in $\mathcal L(Y)$; otherwise the
full alphabet bound applies. Rearrangement proves the claim.
The event may depend on both $X$ and $Y$; no independence is used.
\end{proof}
